% ============================================================
%  File: Appendix1.tex (Linear Algebra)
%  \input this file as an appendix chapter.
%  All environments reference the existing preamble.
%  British English throughout.
% ============================================================

\chapter{Linear Algebra: A Self-Study Reference}
\label{app:linalg}

This appendix provides a concise but complete review of the linear
algebra concepts used throughout these lecture notes.  Each topic is
presented with a formal definition, worked examples by hand, and the
corresponding MATLAB commands so that you can verify every result
interactively.  No prior knowledge beyond basic school-level algebra is
assumed.

All MATLAB code in this appendix can be typed directly into the MATLAB
Command Window or saved as a script.  Matrices are entered row by row,
with semicolons separating rows and commas (or spaces) separating
elements within a row.

% =============================================================
\section{Scalars, Vectors, and Matrices}
\label{sec:app_basics}
% =============================================================

\begin{defnbox}{Scalar}
	A \textbf{scalar} is a single real (or complex) number, denoted by a
	lower-case italic letter, e.g.\ $a = 3.7$.
\end{defnbox}

\begin{defnbox}{Vector}
	A \textbf{column vector} of size $n$ is an ordered list of $n$ scalars
	arranged vertically:
	\[
	{x} =
	\begin{bmatrix} x_1 \\ x_2 \\ \vdots \\ x_n \end{bmatrix}
	\in \mathbb{R}^n.
	\]
	A {row vector} is the transpose of a column vector:
	${x}^\top = [x_1,\; x_2,\; \ldots,\; x_n]$.  Unless stated
	otherwise, \emph{vector} means column vector in these notes.
\end{defnbox}

\begin{defnbox}{Matrix}
	An $m \times n$ \textbf{matrix} has $m$ rows and $n$ columns:
	\[
	A =
	\begin{bmatrix}
		a_{11} & a_{12} & \cdots & a_{1n} \\
		a_{21} & a_{22} & \cdots & a_{2n} \\
		\vdots & \vdots & \ddots & \vdots \\
		a_{m1} & a_{m2} & \cdots & a_{mn}
	\end{bmatrix}
	\in \mathbb{R}^{m \times n}.
	\]
	The element in row $i$ and column $j$ is written $a_{ij}$ or
	$[A]_{ij}$.  A \textbf{square matrix} has $m = n$.
\end{defnbox}

\begin{notebox}{MATLAB notation}
	MATLAB uses \texttt{(i,j)} indexing for elements, starting from 1:
	\texttt{A(i,j)} returns $a_{ij}$.  Size is queried with
	\texttt{size(A)}, which returns \texttt{[m n]}.
\end{notebox}

\begin{examplebox}{Entering matrices and vectors in MATLAB}
	\begin{lstlisting}[language=Matlab]
		% Column vector
		x = [2; -1; 4];       % 3x1, semicolons separate rows
		
		% Row vector
		y = [5, 0, -3];       % 1x3, commas separate elements
		
		% 3x3 matrix
		A = [1, 2, 3;
		4, 5, 6;
		7, 8, 9];
		
		% Access element in row 2, column 3
		a23 = A(2,3)          % returns 6
		
		% Query dimensions
		[m, n] = size(A)      % m=3, n=3
	\end{lstlisting}
	
	\solution
	Running the code above sets \texttt{a23 = 6} and confirms $A$ is
	$3\times 3$.  Vectors are simply matrices where one dimension equals 1.
\end{examplebox}

% =============================================================
\section{Special Matrices}
\label{sec:app_special}
% =============================================================

Several matrix types appear repeatedly in MPC formulations and are
worth knowing by name.

\begin{defnbox}{Important Special Matrices}
	\begin{itemize}[noitemsep]
		\item \textbf{Zero matrix} $0_{m\times n}$: all entries are
		zero.
		\item \textbf{Identity matrix} $I_n$: square, ones on the main
		diagonal, zeros elsewhere; $[I_n]_{ij} = 1$ if $i=j$, else $0$.
		\item \textbf{Diagonal matrix}: square, all off-diagonal entries are
		zero; written $\mathrm{diag}(d_1, d_2, \ldots, d_n)$.
		\item \textbf{Symmetric matrix}: square, $A = A^\top$, i.e.\
		$a_{ij} = a_{ji}$ for all $i,j$.
		\item \textbf{Positive definite matrix}: symmetric $A$ such that
		${x}^\top A {x} > 0$ for all
		${x} \neq 0$.  Written $A \succ 0$.
		\item \textbf{Positive semi-definite matrix}: symmetric $A$ such that
		${x}^\top A {x} \geq 0$.  Written $A \succeq 0$.
	\end{itemize}
\end{defnbox}

\begin{notebox}{Why positive definiteness matters in MPC}
	The cost matrices $Q$ and $R$ in the MPC objective must satisfy
	$Q \succeq 0$ and $R \succ 0$ for the optimisation problem to be
	well-posed (bounded below with a unique minimiser).
\end{notebox}

\begin{examplebox}{Special matrices in MATLAB}
	\begin{lstlisting}[language=Matlab]
		n = 4;
		Z  = zeros(n);          % 4x4 zero matrix
		I  = eye(n);            % 4x4 identity matrix
		D  = diag([1, 2, 3, 4]);% 4x4 diagonal matrix
		
		% Check symmetry
		A  = [4, 2; 2, 3];
		issymmetric(A)          % returns 1 (true)
		
		% Check positive definiteness via eigenvalues
		ev = eig(A)             % all positive => A is positive definite
	\end{lstlisting}
	\solution
	\texttt{eig(A)} returns \texttt{[1.4384; 5.5616]}, both positive,
	confirming $A \succ 0$.
\end{examplebox}

% =============================================================
\section{Matrix Addition and Subtraction}
\label{sec:app_addsubtract}
% =============================================================

\begin{defnbox}{Addition and Subtraction}
	Two matrices $A, B \in \mathbb{R}^{m \times n}$ of the \emph{same
		size} are added (or subtracted) element by element:
	\[
	[A \pm B]_{ij} = a_{ij} \pm b_{ij}.
	\]
	Addition with matrices of different sizes is \textbf{not defined}.
\end{defnbox}

\textbf{Properties of matrix addition:}
\begin{itemize}[noitemsep]
	\item Commutativity: $A + B = B + A$
	\item Associativity: $(A + B) + C = A + (B + C)$
	\item Identity element: $A + 0 = A$
	\item Distributivity over scalars: $\alpha(A+B) = \alpha A + \alpha B$
\end{itemize}

\begin{examplebox}{Matrix addition and subtraction by hand and in MATLAB}
	Let
	\[
	A = \begin{bmatrix} 1 & 3 \\ -2 & 4 \end{bmatrix},
	\qquad
	B = \begin{bmatrix} 5 & -1 \\ 0 & 2 \end{bmatrix}.
	\]
	Compute $A + B$, $A - B$, and $3A$.
	
	\solution
	\[
	A + B =
	\begin{bmatrix} 1+5 & 3-1 \\ -2+0 & 4+2 \end{bmatrix}
	= \begin{bmatrix} 6 & 2 \\ -2 & 6 \end{bmatrix},
	\]
	\[
	A - B =
	\begin{bmatrix} 1-5 & 3+1 \\ -2-0 & 4-2 \end{bmatrix}
	= \begin{bmatrix} -4 & 4 \\ -2 & 2 \end{bmatrix},
	\]
	\[
	3A =
	\begin{bmatrix} 3 & 9 \\ -6 & 12 \end{bmatrix}.
	\]
	
	\begin{lstlisting}[language=Matlab]
		A = [1, 3; -2, 4];
		B = [5, -1; 0, 2];
		
		ApB = A + B     % [6  2; -2  6]
		AmB = A - B     % [-4  4; -2  2]
		threeA = 3*A    % [3  9; -6  12]
	\end{lstlisting}
\end{examplebox}

% =============================================================
\section{Matrix Transpose}
\label{sec:app_transpose}
% =============================================================

\begin{defnbox}{Transpose}
	The \textbf{transpose} of $A \in \mathbb{R}^{m \times n}$ is the
	matrix $A^\top \in \mathbb{R}^{n \times m}$ obtained by swapping rows
	and columns: $[A^\top]_{ij} = a_{ji}$.
\end{defnbox}

\textbf{Useful transpose identities:}
\begin{align}
	(A^\top)^\top &= A \label{eq:t1}\\
	(A + B)^\top  &= A^\top + B^\top \label{eq:t2}\\
	(AB)^\top     &= B^\top A^\top \label{eq:t3}\\
	(\alpha A)^\top &= \alpha A^\top \label{eq:t4}
\end{align}

Identity \eqref{eq:t3} is the most frequently used and most easily
forgotten: note the \emph{reversed} order.

\begin{examplebox}{Transpose in MATLAB}
	\begin{lstlisting}[language=Matlab]
		A  = [1, 2, 3; 4, 5, 6];   % 2x3 matrix
		At = A'                      % 3x2 transpose
		
		% Verify (AB)' = B'A'
		B  = [1, 0; -1, 2; 3, 1];  % 3x2 so that AB is defined
		lhs = (A*B)'
		rhs = B' * A'
		isequal(lhs, rhs)           % returns 1 (true)
	\end{lstlisting}
	\solution
	\texttt{At} is the $3\times2$ matrix with rows \texttt{[1 4]},
	\texttt{[2 5]}, \texttt{[3 6]}.  MATLAB confirms the identity.
	
	\textbf{Warning:} in MATLAB, \texttt{A'} computes the
	\emph{conjugate} (Hermitian) transpose.  For real matrices this makes
	no difference.  For complex matrices use \texttt{A.'} (dot-transpose)
	if you want the non-conjugate transpose.
\end{examplebox}

% =============================================================
\section{Matrix--Matrix Multiplication}
\label{sec:app_matmat}
% =============================================================

\begin{defnbox}{Matrix Multiplication}
	If $A \in \mathbb{R}^{m \times p}$ and $B \in \mathbb{R}^{p \times n}$
	(the inner dimensions must match), their product
	$C = AB \in \mathbb{R}^{m \times n}$ is defined by
	\[
	c_{ij} = \sum_{k=1}^{p} a_{ik}\, b_{kj},
	\qquad i = 1,\ldots,m,\quad j = 1,\ldots,n.
	\]
	In words: the $(i,j)$ entry of $C$ is the \emph{dot product} of the
	$i$-th row of $A$ with the $j$-th column of $B$.
\end{defnbox}

\begin{alertbox}{Matrix multiplication is NOT commutative}
	In general $AB \neq BA$.  The product $BA$ may not even be defined if
	$A$ and $B$ are not square.  Always check dimension compatibility
	before multiplying: an $m\times p$ matrix times a $p \times n$ matrix
	gives an $m \times n$ result.
\end{alertbox}

\textbf{Properties:}
\begin{itemize}[noitemsep]
	\item Associativity: $(AB)C = A(BC)$
	\item Distributivity: $A(B+C) = AB + AC$
	\item Identity: $AI = IA = A$
	\item Scalar: $(\alpha A)B = A(\alpha B) = \alpha(AB)$
\end{itemize}

\begin{examplebox}{Matrix--matrix multiplication by hand and in MATLAB}
	Let
	\[
	A = \begin{bmatrix} 1 & 2 \\ 3 & 4 \end{bmatrix},
	\qquad
	B = \begin{bmatrix} 5 & 6 \\ 7 & 8 \end{bmatrix}.
	\]
	Compute $AB$ and $BA$ and confirm they differ.
	
	\solution
	\[
	AB =
	\begin{bmatrix}
		1\cdot5 + 2\cdot7 & 1\cdot6 + 2\cdot8 \\
		3\cdot5 + 4\cdot7 & 3\cdot6 + 4\cdot8
	\end{bmatrix}
	=
	\begin{bmatrix} 19 & 22 \\ 43 & 50 \end{bmatrix},
	\]
	\[
	BA =
	\begin{bmatrix}
		5\cdot1 + 6\cdot3 & 5\cdot2 + 6\cdot4 \\
		7\cdot1 + 8\cdot3 & 7\cdot2 + 8\cdot4
	\end{bmatrix}
	=
	\begin{bmatrix} 23 & 34 \\ 31 & 46 \end{bmatrix}.
	\]
	Clearly $AB \neq BA$.
	
	\begin{lstlisting}[language=Matlab]
		A  = [1, 2; 3, 4];
		B  = [5, 6; 7, 8];
		
		AB = A * B          % [19 22; 43 50]
		BA = B * A          % [23 34; 31 46]
		isequal(AB, BA)     % returns 0 (false) -- not commutative
	\end{lstlisting}
\end{examplebox}

\begin{examplebox}{Non-square matrix multiplication}
	\[
	A = \begin{bmatrix} 1 & 0 & -1 \\ 2 & 3 & 1 \end{bmatrix} \in
	\mathbb{R}^{2\times 3}, \qquad
	B = \begin{bmatrix} 2 \\ -1 \\ 4 \end{bmatrix} \in \mathbb{R}^{3\times 1}.
	\]
	Compute $AB$ (which exists) and state why $BA$ does not exist.
	
	\solution
	$A$ is $2\times 3$ and $B$ is $3\times 1$: inner dimensions both equal
	3, so $AB \in \mathbb{R}^{2\times 1}$:
	\[
	AB =
	\begin{bmatrix}
		1\cdot2 + 0\cdot(-1) + (-1)\cdot4 \\
		2\cdot2 + 3\cdot(-1) + 1\cdot4
	\end{bmatrix}
	=
	\begin{bmatrix} -2 \\ 5 \end{bmatrix}.
	\]
	$BA$ would require $B$ ($3\times1$) times $A$ ($2\times3$): inner
	dimensions are $1 \neq 2$, so $BA$ is \textbf{undefined}.
	
	\begin{lstlisting}[language=Matlab]
		A = [1, 0, -1; 2, 3, 1];
		B = [2; -1; 4];
		AB = A * B          % [-2; 5]
		BA = B * A          % Error: inner matrix dimensions must agree
	\end{lstlisting}
\end{examplebox}

% =============================================================
\section{Matrix--Vector Multiplication}
\label{sec:app_matvec}
% =============================================================

Matrix--vector multiplication is a special case of matrix--matrix
multiplication where $n=1$, but it has a particularly important
geometric interpretation: $A{x}$ transforms the vector
${x}$ into a new vector in a (possibly different) space.

\begin{defnbox}{Matrix--Vector Product}
	For $A \in \mathbb{R}^{m\times n}$ and
	${x} \in \mathbb{R}^n$:
	\[
	{y} = A{x}, \qquad
	y_i = \sum_{j=1}^{n} a_{ij}\, x_j, \quad i = 1,\ldots,m.
	\]
	Equivalently, $A{x}$ is a {linear combination} of the
	columns of $A$ with coefficients given by the entries of ${x}$:
	\[
	A{x} = x_1 {a}_1 + x_2 {a}_2 + \cdots
	+ x_n {a}_n,
	\]
	where ${a}_j$ denotes the $j$-th column of $A$.
\end{defnbox}

In state-space control the plant update $x_{k+1} = Ax_k + Bu_k$ is
nothing more than two matrix--vector products added together.

\begin{examplebox}{Matrix--vector multiplication by hand and in MATLAB}
	Given the state-space matrices
	\[
	A = \begin{bmatrix} 0.9 & 0.1 \\ 0 & 0.8 \end{bmatrix},
	\quad
	B = \begin{bmatrix} 0.1 \\ 0.5 \end{bmatrix},
	\quad
	x_k = \begin{bmatrix} 2 \\ -1 \end{bmatrix},
	\quad
	u_k = 3,
	\]
	compute the next state $x_{k+1} = Ax_k + Bu_k$.
	
	\solution
	\[
	Ax_k =
	\begin{bmatrix} 0.9\cdot2 + 0.1\cdot(-1) \\ 0\cdot2 + 0.8\cdot(-1)
	\end{bmatrix}
	= \begin{bmatrix} 1.7 \\ -0.8 \end{bmatrix},
	\qquad
	Bu_k =
	\begin{bmatrix} 0.3 \\ 1.5 \end{bmatrix},
	\]
	\[
	x_{k+1} = \begin{bmatrix} 1.7 \\ -0.8 \end{bmatrix}
	+ \begin{bmatrix} 0.3 \\ 1.5 \end{bmatrix}
	= \begin{bmatrix} 2.0 \\ 0.7 \end{bmatrix}.
	\]
	
	\begin{lstlisting}[language=Matlab]
		A  = [0.9, 0.1; 0, 0.8];
		B  = [0.1; 0.5];
		xk = [2; -1];
		uk = 3;
		
		xk1 = A*xk + B*uk    % [2.0; 0.7]
	\end{lstlisting}
\end{examplebox}

% =============================================================
\section{The Inverse of a Matrix}
\label{sec:app_inverse}
% =============================================================

\begin{defnbox}{Matrix Inverse}
	A square matrix $A \in \mathbb{R}^{n\times n}$ is \textbf{invertible}
	(or \textbf{non-singular}) if there exists a matrix
	$A^{-1} \in \mathbb{R}^{n\times n}$ such that
	\[
	A A^{-1} = A^{-1} A = I_n.
	\]
	If no such matrix exists, $A$ is called \textbf{singular}.
\end{defnbox}

\begin{defnbox}{Conditions for Invertibility}
	A square matrix $A$ is invertible if and only if any one of the
	following equivalent conditions holds:
	\begin{itemize}[noitemsep]
		\item $\det(A) \neq 0$ (see Section~\ref{sec:app_det});
		\item all eigenvalues of $A$ are non-zero
		(see Section~\ref{sec:app_eig});
		\item the columns (or rows) of $A$ are linearly independent;
		\item the equation $A{x} = 0$ has only the trivial
		solution ${x} = 0$.
	\end{itemize}
\end{defnbox}

\textbf{Useful inverse identities:}
\begin{align}
	(A^{-1})^{-1} &= A \\
	(AB)^{-1}     &= B^{-1} A^{-1} \quad \text{(reversed order)} \\
	(A^\top)^{-1} &= (A^{-1})^\top \\
	(\alpha A)^{-1} &= \tfrac{1}{\alpha} A^{-1}
\end{align}

\textbf{Inverse of a $2\times2$ matrix} by hand:
\[
A = \begin{bmatrix} a & b \\ c & d \end{bmatrix}
\implies
A^{-1} = \frac{1}{ad - bc}
\begin{bmatrix} d & -b \\ -c & a \end{bmatrix}.
\]

\begin{examplebox}{Computing the inverse by hand and in MATLAB}
	Let
	\[
	A = \begin{bmatrix} 3 & 1 \\ 5 & 2 \end{bmatrix}.
	\]
	Find $A^{-1}$ and verify $AA^{-1} = I$.
	
	\solution
	$\det(A) = 3\cdot2 - 1\cdot5 = 1$, so
	\[
	A^{-1} = \frac{1}{1}
	\begin{bmatrix} 2 & -1 \\ -5 & 3 \end{bmatrix}
	= \begin{bmatrix} 2 & -1 \\ -5 & 3 \end{bmatrix}.
	\]
	Verification:
	\[
	AA^{-1} =
	\begin{bmatrix} 3 & 1 \\ 5 & 2 \end{bmatrix}
	\begin{bmatrix} 2 & -1 \\ -5 & 3 \end{bmatrix}
	=
	\begin{bmatrix} 6-5 & -3+3 \\ 10-10 & -5+6 \end{bmatrix}
	= \begin{bmatrix} 1 & 0 \\ 0 & 1 \end{bmatrix} = I. \checkmark
	\]
	
	\begin{lstlisting}[language=Matlab]
		A    = [3, 1; 5, 2];
		Ainv = inv(A)           % [2 -1; -5 3]
		
		% Verify
		A * Ainv                % should give eye(2)
		norm(A*Ainv - eye(2))   % numerical error, should be ~1e-15
	\end{lstlisting}
\end{examplebox}

\begin{alertbox}{Avoid \texttt{inv(A)} in numerical computations}
	Although \texttt{inv(A)} is convenient for checking results,
	numerically it is almost always better to use the backslash operator
	\texttt{A\textbackslash b} to solve $Ax = b$
	directly, without forming $A^{-1}$ explicitly.  The backslash
	operator uses LU factorisation, which is faster and more numerically
	stable than computing the full inverse.
	
	\begin{lstlisting}[language=Matlab]
		% Solve Ax = b
		b = [1; 2];
		x_inv = inv(A) * b    % slow and less accurate
		x_bs  = A \ b         % fast, numerically recommended
	\end{lstlisting}
\end{alertbox}

\begin{examplebox}{Detecting a singular matrix}
	\begin{lstlisting}[language=Matlab]
		A_sing = [1, 2; 2, 4];     % row 2 = 2 * row 1
		
		det(A_sing)                 % returns 0 => singular
		rank(A_sing)                % returns 1 (not full rank)
		cond(A_sing)                % returns Inf
		
		inv(A_sing)                 % MATLAB warns: matrix is singular
	\end{lstlisting}
	\solution
	The determinant is zero and the rank is 1 (not 2), confirming
	singularity.  MATLAB returns a matrix of \texttt{Inf} values from
	\texttt{inv} and issues a warning.  The condition number
	\texttt{cond(A)} measures how close a matrix is to being singular;
	values much larger than $10^6$ indicate near-singularity.
\end{examplebox}

% =============================================================
\section{The Determinant}
\label{sec:app_det}
% =============================================================

\begin{defnbox}{Determinant}
	The \textbf{determinant} $\det(A)$ (also written $|A|$) is a scalar
	function of a square matrix that encodes fundamental properties such as
	invertibility, volume scaling, and orientation.
	
	For a $2\times2$ matrix:
	\[
	\det\!\begin{bmatrix} a & b \\ c & d \end{bmatrix} = ad - bc.
	\]
	For a $3\times3$ matrix (cofactor expansion along the first row):
	\[
	\det\!\begin{bmatrix} a_{11} & a_{12} & a_{13} \\
		a_{21} & a_{22} & a_{23} \\
		a_{31} & a_{32} & a_{33} \end{bmatrix}
	= a_{11}(a_{22}a_{33} - a_{23}a_{32})
	- a_{12}(a_{21}a_{33} - a_{23}a_{31})
	+ a_{13}(a_{21}a_{32} - a_{22}a_{31}).
	\]
\end{defnbox}

\textbf{Key properties:}
\begin{itemize}[noitemsep]
	\item $\det(I) = 1$
	\item $\det(AB) = \det(A)\det(B)$
	\item $\det(A^\top) = \det(A)$
	\item $\det(A^{-1}) = 1/\det(A)$
	\item $\det(\alpha A) = \alpha^n \det(A)$ for $A \in \mathbb{R}^{n\times n}$
	\item $\det(A) = 0$ if and only if $A$ is singular
\end{itemize}

\begin{examplebox}{Determinant by hand and in MATLAB}
	Compute $\det(A)$ for
	\[
	A = \begin{bmatrix} 2 & -1 & 0 \\ 1 & 3 & -2 \\ 0 & 1 & 4 \end{bmatrix}.
	\]
	
	\solution
	Expanding along the first row:
	\[
	\det(A) = 2\bigl(3\cdot4 - (-2)\cdot1\bigr)
	- (-1)\bigl(1\cdot4 - (-2)\cdot0\bigr)
	+ 0\bigl(\cdots\bigr)
	= 2(14) + 1(4) = 32.
	\]
	
	\begin{lstlisting}[language=Matlab]
		A    = [2, -1, 0; 1, 3, -2; 0, 1, 4];
		d    = det(A)           % returns 32
	\end{lstlisting}
\end{examplebox}

% =============================================================
\section{The Trace}
\label{sec:app_trace}
% =============================================================

\begin{defnbox}{Trace}
	The \textbf{trace} of a square matrix $A \in \mathbb{R}^{n\times n}$
	is the sum of its diagonal elements:
	\[
	\mathrm{tr}(A) = \sum_{i=1}^{n} a_{ii} = a_{11} + a_{22} + \cdots + a_{nn}.
	\]
\end{defnbox}

\textbf{Key properties:}
\begin{itemize}[noitemsep]
	\item Linearity: $\mathrm{tr}(\alpha A + \beta B) = \alpha\,\mathrm{tr}(A) + \beta\,\mathrm{tr}(B)$
	\item Cyclic permutation: $\mathrm{tr}(ABC) = \mathrm{tr}(CAB) = \mathrm{tr}(BCA)$
	\item Transpose invariance: $\mathrm{tr}(A) = \mathrm{tr}(A^\top)$
	\item Equal to sum of eigenvalues: $\mathrm{tr}(A) = \sum_{i=1}^n \lambda_i$
\end{itemize}

The trace appears in control theory when expressing quadratic cost
functions compactly.  The MPC stage cost
${x}^\top Q {x}$ can be written as
$\mathrm{tr}(Q{x}{x}^\top)$, which is useful in
stochastic and robust formulations.

\begin{examplebox}{Trace in MATLAB}
	\begin{lstlisting}[language=Matlab]
		A = [1, 2, 3;
		4, 5, 6;
		7, 8, 9];
		
		tr = trace(A)           % 1 + 5 + 9 = 15
		
		% Verify trace equals sum of eigenvalues
		lambda = eig(A);
		sum(lambda)             % also 15 (up to numerical precision)
		
		% Cyclic property: tr(ABC) = tr(CAB)
		B = rand(3);  C = rand(3);
		abs(trace(A*B*C) - trace(C*A*B))   % ~1e-14, effectively zero
	\end{lstlisting}
	\solution
	\texttt{trace(A) = 15 = 1+5+9}.  The sum of eigenvalues agrees
	(numerically), confirming the identity.
\end{examplebox}

% =============================================================
\section{Eigenvalues and Eigenvectors}
\label{sec:app_eig}
% =============================================================

Eigenvalues and eigenvectors are central to understanding the dynamic
behaviour of linear systems.  In discrete-time state-space control,
the eigenvalues of the system matrix $A$ determine stability: the
closed-loop system $x_{k+1} = A_\mathrm{cl} x_k$ is stable if and
only if all eigenvalues of $A_\mathrm{cl}$ lie strictly inside the
unit circle.

\begin{defnbox}{Eigenvalue and Eigenvector}
	Let $A \in \mathbb{R}^{n\times n}$.  A scalar $\lambda \in \mathbb{C}$
	is an \textbf{eigenvalue} of $A$ if there exists a non-zero vector
	${v} \in \mathbb{C}^n$ such that
	\[
	A{v} = \lambda {v}.
	\]
	The vector $v$ is the corresponding \textbf{eigenvector}.
	Geometrically, $A$ does not rotate ${v}$; it only scales it by
	$\lambda$.
\end{defnbox}

\textbf{Finding eigenvalues — the characteristic equation:}

Rearranging $A{v} = \lambda{v}$ gives
$(A - \lambda I){v} = {0}$.  For a non-trivial solution
${v} \neq {0}$, the matrix $(A-\lambda I)$ must be
singular:
\[
\det(A - \lambda I) = 0.
\]
This is the \textbf{characteristic equation}, a degree-$n$ polynomial
in $\lambda$ whose $n$ roots (counting multiplicity, possibly complex)
are the eigenvalues.

\begin{examplebox}{Eigenvalues by hand for a $2\times2$ matrix}
	Find the eigenvalues of
	\[
	A = \begin{bmatrix} 3 & 1 \\ 1 & 3 \end{bmatrix}.
	\]
	
	\solution
	\[
	\det(A - \lambda I)
	= \det\begin{bmatrix} 3-\lambda & 1 \\ 1 & 3-\lambda \end{bmatrix}
	= (3-\lambda)^2 - 1 = 0.
	\]
	\[
	(3-\lambda)^2 = 1
	\implies 3 - \lambda = \pm 1
	\implies \lambda_1 = 4, \quad \lambda_2 = 2.
	\]
	Both eigenvalues are real and positive.  For a discrete-time system
	matrix, we would check $|\lambda_i| < 1$; here both are $> 1$, so this
	$A$ would represent an unstable open-loop plant.
	
	\begin{lstlisting}[language=Matlab]
		A      = [3, 1; 1, 3];
		lambda = eig(A)         % [2; 4]
	\end{lstlisting}
\end{examplebox}

\begin{examplebox}{Eigenvectors and the full eigendecomposition}
	Using the same matrix $A = \begin{bmatrix} 3 & 1 \\ 1 & 3 \end{bmatrix}$,
	find the eigenvectors corresponding to $\lambda_1 = 4$ and
	$\lambda_2 = 2$.
	
	\solution
	\textbf{For $\lambda_1 = 4$:}
	\[
	(A - 4I)v = \begin{bmatrix} -1 & 1 \\ 1 & -1 \end{bmatrix}
	\begin{bmatrix} v_1 \\ v_2 \end{bmatrix} = 0
	\implies -v_1 + v_2 = 0 \implies v_1 = v_2.
	\]
	Choosing $v_1 = 1$: $v_1 = \tfrac{1}{\sqrt{2}}[1,\;1]^\top$
	(normalised).
	
	\textbf{For $\lambda_2 = 2$:}
	\[
	(A - 2I)v = \begin{bmatrix} 1 & 1 \\ 1 & 1 \end{bmatrix}
	v = 0
	\implies v_1 + v_2 = 0 \implies v_2 = -v_1.
	\]
	Choosing $v_1 = 1$: $v_2 = \tfrac{1}{\sqrt{2}}[1,\;-1]^\top$.
	
	\begin{lstlisting}[language=Matlab]
		A = [3, 1; 1, 3];
		
		% eig returns V (eigenvector matrix) and D (diagonal eigenvalue matrix)
		[V, D] = eig(A)
		% V: columns are eigenvectors
		% D: diagonal entries are eigenvalues
		
		% Verify A*v = lambda*v for first eigenvector
		v1     = V(:,1);
		lam1   = D(1,1);
		norm(A*v1 - lam1*v1)    % ~1e-15, effectively zero
		
		% The eigendecomposition: A = V * D * inv(V)
		norm(A - V*D*inv(V))    % ~1e-15, effectively zero
	\end{lstlisting}
\end{examplebox}

\begin{defnbox}{Eigenvalue Properties for Stability Analysis}
	For a discrete-time system $x_{k+1} = A x_k$:
	\begin{itemize}[noitemsep]
		\item \textbf{Asymptotically stable:} all eigenvalues satisfy
		$|\lambda_i| < 1$ (strictly inside the unit circle).
		\item \textbf{Marginally stable:} all $|\lambda_i| \leq 1$ and
		those on the unit circle have distinct eigenvectors.
		\item \textbf{Unstable:} at least one eigenvalue satisfies
		$|\lambda_i| > 1$.
	\end{itemize}
	For a continuous-time system $\dot{x} = Ax$, replace the unit circle
	with the left half-plane: stable iff $\mathrm{Re}(\lambda_i) < 0$.
\end{defnbox}

\begin{examplebox}{Stability check via eigenvalues in MATLAB}
	Check whether the following discrete-time system matrix is stable:
	\[
	A = \begin{bmatrix} 0.98 & 0.01 \\ 0.02 & 0.97 \end{bmatrix}.
	\]
	
	\solution
	\begin{lstlisting}[language=Matlab]
		A      = [0.98, 0.01; 0.02, 0.97];
		lambda = eig(A)             % [0.96; 0.99]
		abs_lambda = abs(lambda)    % [0.96; 0.99]
		all(abs_lambda < 1)         % returns 1 (true) => stable
		
		% Visualise on the unit circle
		theta = linspace(0, 2*pi, 200);
		figure; hold on;
		plot(cos(theta), sin(theta), 'k--');    % unit circle
		plot(real(lambda), imag(lambda), 'rx', 'MarkerSize', 10);
		axis equal; grid on;
		xlabel('Real'); ylabel('Imaginary');
		title('Eigenvalues of A');
	\end{lstlisting}
	Both eigenvalues have magnitude less than 1, so the system is
	asymptotically stable.  The plot shows both poles inside the unit
	circle.
\end{examplebox}

% =============================================================
\section{Putting It All Together: A State-Space Example}
\label{sec:app_together}
% =============================================================

This section demonstrates how all the operations above arise together
in a single state-space simulation, mirroring the structure of the MPC
plant model used throughout these notes.

\begin{examplebox}{Two-state discrete system: simulation and analysis}
	Consider the discrete-time system
	\[
	x_{k+1} = A x_k + B u_k, \qquad y_k = C x_k,
	\]
	with
	\[
	A = \begin{bmatrix} 0.9 & 0.1 \\ 0.0 & 0.8 \end{bmatrix},
	\quad
	B = \begin{bmatrix} 0.1 \\ 0.3 \end{bmatrix},
	\quad
	C = \begin{bmatrix} 1 & 0 \end{bmatrix},
	\quad
	x_0 = \begin{bmatrix} 5 \\ 2 \end{bmatrix},
	\quad u_k = 0\;\forall k.
	\]
	
	Carry out the following tasks in MATLAB:
	\begin{enumerate}[noitemsep]
		\item Confirm the system is stable.
		\item Compute $\mathrm{tr}(A)$ and $\det(A)$ and relate them to the
		eigenvalues.
		\item Simulate 20 steps and plot the output $y_k$.
	\end{enumerate}
	
	\solution
	\begin{lstlisting}[language=Matlab]
		A  = [0.9, 0.1; 0.0, 0.8];
		B  = [0.1; 0.3];
		C  = [1, 0];
		x  = [5; 2];
		N  = 20;
		
		%% 1. Stability
		lambda = eig(A);
		disp('Eigenvalues:'); disp(lambda)
		disp('Magnitudes:');  disp(abs(lambda))
		% Both < 1 => stable
		
		%% 2. Trace and determinant
		tr_A  = trace(A)          % 0.9 + 0.8 = 1.7
		det_A = det(A)            % 0.9*0.8 - 0.1*0 = 0.72
		sum_lam  = sum(lambda)    % also 1.7 = trace
		prod_lam = prod(lambda)   % also 0.72 = det
		
		%% 3. Simulation
		y = zeros(1, N);
		for k = 1:N
		y(k) = C * x;
		x    = A * x + B * 0;   % u=0
		end
		
		figure;
		stem(1:N, y, 'filled');
		xlabel('Time step k'); ylabel('y_k');
		title('Output y_k (free response, u=0)');
		grid on;
	\end{lstlisting}
	
	The output decays to zero, consistent with the stable eigenvalues
	$\lambda_1 \approx 0.9$ and $\lambda_2 = 0.8$.  Notice that
	$\mathrm{tr}(A) = \sum \lambda_i = 1.7$ and
	$\det(A) = \prod \lambda_i = 0.72$, confirming the identities from
	Sections~\ref{sec:app_trace} and \ref{sec:app_det}.
\end{examplebox}

% =============================================================
\section{Quick Reference: MATLAB Commands}
\label{sec:app_matlab_ref}
% =============================================================

\begin{center}
	\renewcommand{\arraystretch}{1.35}
	\begin{tabular}{>{\ttfamily\small}l >{\small}l >{\small}l}
		\toprule
		\normalfont\textbf{MATLAB command} & \textbf{Operation} &
		\textbf{Example} \\
		\midrule
		A + B               & Addition              & Element-wise, same size \\
		A - B               & Subtraction           & Element-wise, same size \\
		alpha * A           & Scalar multiplication & Multiplies every entry \\
		A * B               & Matrix multiplication & Inner dims must match \\
		A'                  & Transpose (conjugate) & Use \texttt{A.'} for real transpose \\
		A \textasciicircum{} n & Matrix power       & $A^n$, square $A$ only \\
		inv(A)              & Matrix inverse        & Avoid for solving $Ax=b$ \\
		A \textbackslash{} b & Solve $Ax = b$       & Preferred over \texttt{inv(A)*b} \\
		det(A)              & Determinant           & Scalar result \\
		trace(A)            & Trace                 & Sum of diagonal \\
		eig(A)              & Eigenvalues           & Column vector $\lambda$ \\
		{[}V,D{]} = eig(A)  & Eigenvectors + values & Columns of \texttt{V}, diag of \texttt{D} \\
		rank(A)             & Rank                  & Number of independent columns \\
		cond(A)             & Condition number      & Large value $\Rightarrow$ near-singular \\
		norm(A)             & 2-norm (default)      & \texttt{norm(A,1)}, \texttt{norm(A,'fro')} also available \\
		size(A)             & Dimensions            & Returns \texttt{[m n]} \\
		zeros(m,n)          & Zero matrix           & All entries 0 \\
		ones(m,n)           & Ones matrix           & All entries 1 \\
		eye(n)              & Identity matrix       & $I_n$ \\
		diag(v)             & Diagonal matrix       & \texttt{diag(A)} extracts diagonal of $A$ \\
		linspace(a,b,n)     & Equally spaced vector & $n$ points from $a$ to $b$ \\
		\bottomrule
	\end{tabular}
\end{center}

\begin{summarybox}{Linear Algebra Summary for MPC}
	\begin{enumerate}[noitemsep]
		\item \textbf{Dimensions matter:} always check that inner dimensions
		are compatible before multiplying; $A_{m\times p} \cdot
		B_{p\times n}$ gives $C_{m\times n}$.
		\item \textbf{Multiplication is not commutative:} $AB \neq BA$ in
		general; the reversed-order rule $(AB)^\top = B^\top A^\top$
		and $(AB)^{-1} = B^{-1}A^{-1}$ are essential.
		\item \textbf{Invertibility:} $A^{-1}$ exists iff $\det(A)\neq 0$
		iff all eigenvalues are non-zero.  Use \texttt{A\textbackslash
			b} rather than \texttt{inv(A)*b}.
		\item \textbf{Stability:} for discrete-time systems, check
		$|\lambda_i| < 1$ for all eigenvalues of the closed-loop
		matrix.
		\item \textbf{Trace and determinant} equal the sum and product of
		eigenvalues respectively, providing a fast sanity check.
		\item \textbf{Positive definiteness} of cost matrices $Q \succeq 0$,
		$R \succ 0$ is required for the MPC QP to be convex and
		bounded below; verify with \texttt{eig(Q)} and
		\texttt{eig(R)}.
	\end{enumerate}
\end{summarybox}